\documentclass[fleqn]{article}
\oddsidemargin 0.0in
\textwidth 6.0in
\thispagestyle{empty}
\usepackage{import}
\usepackage{amsmath}
\usepackage{bigints}
\usepackage{graphicx}
\usepackage[english]{babel}
\usepackage[utf8x]{inputenc}
\usepackage{float}
\usepackage[colorinlistoftodos]{todonotes}

\definecolor{hwColor}{HTML}{AD53BA}

\begin{document}

  \begin{titlepage}

    \newcommand{\HRule}{\rule{\linewidth}{0.5mm}} % Defines a new command for the horizontal lines, change thickness here

    \center % Center everything on the page

    \textsc{\LARGE Arizona State University}\\[1.5cm] % Name of your university/college

    \textsc{\LARGE Mathematical Methods For Physics I }\\[1.5cm] % Major heading such as course name


    \begin{figure}
      \includegraphics[width=\linewidth]{asu.png}
    \end{figure}


    \HRule \\[0.4cm]
    { \huge \bfseries Homework 14}\\[0.4cm] 
    \HRule \\[1.5cm]

    \textbf{Behnam Amiri}

    \bigbreak

    \textbf{Prof: Cecilia Lunardini}

    \bigbreak


    \textbf{{\large \today}\\[2cm]}

    \vfill % Fill the rest of the page with whitespace

  \end{titlepage}

  \begin{enumerate}
    \item  Solve the problem of the vibrating string of length $L$ (covered in class) with the function
      $$f(x)=\begin{cases} \frac{4 h x}{L} &\mbox{if } 0 \leq x < L/4 \\ 
      2 h \left( 1 - \frac{2x}{L}\right) &\mbox{if }  L/4  \leq x \leq L/2 \\ 
      0 & \mbox{if } x > L/2 \end{cases} $$
      describing the initial displacement of the string at $t=0$. Plot $f(x)$. Assume that initially the string is released from rest. [Note: please do all the steps, including the separation of the variables, even though many of them have been shown in class. ]

      \textcolor{hwColor}{
        $
          \dfrac{\partial^2y}{\partial x^2}-\dfrac{1}{v^2}\dfrac{\partial^y}{\partial t^2}=0
          \\
          \\
          y(x,t)=X(x)T(t) \\
          \\
          \dfrac{\partial^2}{\partial x^2}\left(X(x)T(t)\right)-\dfrac{1}{v^2}\dfrac{\partial^2}{\partial t^2}\left(X(x)T(t)\right)=0 \\
          \\
          T(t)\dfrac{d^2X}{dx^2}-\dfrac{1}{v^2}X(x)\dfrac{d^2T}{dt^2}=0 \\
          \\
        $ 
        Now divide both sides by $\dfrac{1}{X(x)T(t)}$, so we get:
        $
          \dfrac{1}{X(x)}\dfrac{d^2X}{dx^2}-\dfrac{1}{v^2}\dfrac{1}{T(t)}\dfrac{d^2T}{dt^2}=0 \\
          \\
          \\
          \dfrac{1}{X}\dfrac{d^2x}{dx^2}=\dfrac{1}{v^2}\dfrac{1}{T}\dfrac{d^2T}{dt^2} \Rightarrow \begin{cases}
            X(x)=X \\
            T(t)=T
          \end{cases}
        $ \\
        \\
        \\
        $
          \begin{cases}
            \dfrac{d^2X}{dx^2}=-K^2X \\
            \\
            \dfrac{d^2T}{dt^2}=-(vK)^2T
          \end{cases} \Rightarrow \begin{cases}
            \dfrac{d^2X}{dx^2}+K^2X=0 \\
            \\
            \dfrac{d^2T}{dt^2}+(vK)^2T=0
          \end{cases}
        $ where K is a real number. \\
        \\
      }

      \textcolor{hwColor}{
        These two objects on each side of the equation have nothing to do with each other because one of them is a function of x only depends on x. And the other part this on this side only depends on T. 
        How can these two sides be equal? Only when they are constant. \\
        \\
        $
          \begin{cases}
            X(x)=Acos(Kx)+Bsin(Kx) \\
            T(t)=Ccos(Kvt)+Dsin(Kvt)
          \end{cases}
        $
        So for our general solution we have: \\
        \\
        \\
        $
          y(x,t)=\left[Acos(Kx)+Bsin(Kx)\right] \times \left[Ccos(Kvt)+Dsin(Kvt)\right]
        $
        \\
      }

      \textcolor{hwColor}{
        We assume that the ends of the string are fixed. \\
        $
          \begin{cases}
            y(0,t)=0 \Rightarrow (i) ~~ X(0)=0\\
            y(L,t)=0 \Rightarrow (ii) ~~ X(L)=0
          \end{cases}
          \\
          \\
          \\
          (i) ~~ X(x)=Acos(Kx)+Bsin(Kx) \\
          x(0)=Acos(0)+Bsin(0) \longrightarrow A=0 \\
          \\
          \\
          (ii) ~~ X(x)=Bsin(Kx) \\
          X(L)=0 \rightarrow Bsin(KL)=0 \Rightarrow KL=n\pi ~~
        $ where $n=1,2,3,...$ \\
      }

      \textcolor{hwColor}{
        $
          y(x,t)=\sum\limits_{n=1}^{\infty} sin(\dfrac{n\pi}{L}x) \times \left[C_ncos(K_nvt)+D_nsin(K_nvt)\right]
        $
         If $\omega_n=K_nv$ then: \\ 
        $
          \\
          y(x,t)=\sum\limits_{n=1}^{\infty} sin(K_nx)\left[C_ncos(\omega_n t)+D_nsin(\omega_n t)\right]
        $ Where $K_n=\dfrac{n\pi}{L}$
      }

      \textcolor{hwColor}{
        The initial condition is that the string is released from rest which means $y^{'}(x,0)=0$ \\
        \\
        $
          \dfrac{dy}{dt}=\dfrac{d}{dt}\left[\sum\limits_{n=1}^{\infty} sin(K_nx)\left[C_ncos(\omega_n t)+D_nsin(\omega_n t)\right]\right] \\
          \\
          \dfrac{dy}{dt}=\sum\limits_{n=1}^{\infty} \omega_n sin(K_nx)\left[-C_n sin(0)+D_n cos(0)\right]=0 \\
          \\
          \dfrac{dy}{dt}=\sum\limits_{n=1}^{\infty} D_n \omega_n sin(K_nx)=0 \rightarrow D_n=0 \\
          \\
          y(x,t)=\sum\limits_{n=1}^{\infty} C_n sin(K_nx)cos(\omega_n t) \\
          \\
          \\
          y(x,0)=\sum\limits_{n=1}^{\infty} C_n sin(K_nx)cos(0)=f(x) \\
          \\
          y(x,0)=\sum\limits_{n=1}^{\infty} C_n sin(\dfrac{n\pi}{L}x)=f(x) \\
        $
      }

      \textcolor{hwColor}{
        We can find the value of $C_n$ by using the equation 12.6 from the textbook: \\
        $
          C_n=\dfrac{2}{L} \bigints_{0}^{L} f(x) sin(\dfrac{n \pi x}{L})dx \\
          \\
          =\dfrac{2}{L} \left[\bigints_{0}^{\dfrac{L}{4}}\dfrac{4hx}{L}sin(\dfrac{n \pi x}{\dfrac{L}{4}})dx+\bigints_{\dfrac{L}{4}}^{\dfrac{L}{2}}2h(1-\dfrac{2x}{L})sin(\dfrac{n \pi x}{\dfrac{L}{4}})dx+\bigints_{\dfrac{L}{2}}^{L} 0 sin(\dfrac{n \pi x}{\dfrac{L}{4}})dx\right] \\
          \\
          =\dfrac{2}{L} \left[\bigints_{0}^{\dfrac{L}{4}}\dfrac{4hx}{L}sin(\dfrac{n \pi x}{\dfrac{L}{4}})dx+\bigints_{\dfrac{L}{4}}^{\dfrac{L}{2}}2h(1-\dfrac{2x}{L})sin(\dfrac{n \pi x}{\dfrac{L}{4}})dx\right] \\
          \\
          \\
          \\
          \bigints_{0}^{\dfrac{L}{4}}\dfrac{4hx}{L}sin(\dfrac{n \pi x}{\dfrac{L}{4}})dx=\dfrac{4h}{L} \bigints_{0}^{\dfrac{L}{4}}x sin(\dfrac{4\pi nx}{L})dx \\
          \\
          =\dfrac{4h}{L} \left[-\dfrac{Lx}{4 \pi n}cos(\dfrac{4 \pi nx}{L})+\bigints \dfrac{L}{4\pi n} cos(\dfrac{4 \pi nx}{L})dx\right] \\
          \\
          =\dfrac{4h}{L} \left[-\dfrac{Lx}{4 \pi n} cos(\dfrac{4 \pi n x}{L})+(\dfrac{L}{4 \pi n})^2 sin(\dfrac{4\pi n x}{L})\right]\Biggr|_{0}^{\dfrac{L}{4}} \\
          \\
          \\
          =\dfrac{4h}{L} \left[-\dfrac{L^2}{16 \pi n}cos(\pi n)+(\dfrac{L}{4 \pi n})^2 sin(\pi n)\right]=-\dfrac{Lh}{4 \pi n}(-1)^n \\
          \\
          \\
          \\
          \bigints_{\dfrac{L}{4}}^{\dfrac{L}{2}}2h(1-\dfrac{2x}{L})sin(\dfrac{n \pi x}{\dfrac{L}{4}})dx=2h\left[\bigints sin(\dfrac{4\pi nx}{L})dx-\bigints \dfrac{2x}{L}sin(\dfrac{4\pi nx}{L})\right]\Biggr|_{\dfrac{L}{4}}^{\dfrac{L}{2}} \\
          \\
          =2h \left[-\dfrac{L}{4 \pi n} cos(\dfrac{4 \pi nx}{L})+\dfrac{x}{2\pi n}cos(\dfrac{4 \pi nx}{L})-\dfrac{L}{8(n\pi)^2}sin(\dfrac{4\pi nx}{L})\right]\Biggr|_{\dfrac{L}{4}}^{\dfrac{L}{2}} \\
          \\
          =2h \left[-\dfrac{L}{4 \pi n}+\dfrac{L}{4 \pi n}+\dfrac{L}{4 \pi n}(-1)^n-\dfrac{1}{8 \pi n}(-1)^n\right]=2h \left[\dfrac{L}{4 \pi n}(-1)^n-\dfrac{1}{8 \pi n}(-1)^n\right] \\
          \\
          =\dfrac{h(-1)^n (2L-1)}{4 \pi n} \\
          \\
          \\
          \\
          \\
          C_n=\dfrac{2}{L} \left[-\dfrac{Lh}{4 \pi n}(-1)^n+\dfrac{h(-1)^n (2L-1)}{4 \pi n}\right] \\
          \\
          \\
          \Longrightarrow C_n=\dfrac{h(-1)^n(L-2L^2)}{4 \pi n} \\
          \\
          \\
          \Longrightarrow y(x,0)=\sum\limits_{n=1}^{\infty} \dfrac{h(-1)^n(L-2L^2)}{4 \pi n} sin(\dfrac{n\pi}{L}x)=f(x) \\
          \\
          \\
          \\
          \\
          \\
          \\
        $
      }

    \item  

      \begin{enumerate}
      \item Derive the differential equation for a  vibrating string in presence of gravity. Assume that the string has constant mass per unit length, $\mu$, and that at rest the string lies horizontally (like in a piano); assume small displacements, purely in the transverse direction (like in the case discussed in class).  

        \textcolor{hwColor}{
          $\tau :$ Tension ~~~~ $\mu:$ Mass per unit length \\
          \\
          $
            F_y=\tau_2 sin(\theta_2)-\tau_1 sin(\theta_1)-mg=\tau_2 cos(\theta_2) \dfrac{sin(\theta_2)}{cos(\theta_2)}-\tau_1 cos(\theta_1) \dfrac{sin(\theta_1)}{cos(\theta_1)}-mg \\
            \\
            F_y=\tau_2 cos(\theta_2) tan(\theta_2)-\tau_1 cos(\theta_1) tan(\theta_1)-mg \\
            \\
          $
          Based on our assumption that all small displacements, purely in the transverse direction we can write the following: \\
          \\
          $
            F_x=\tau_2 cos(\theta_2)-\tau_1 cos(\theta_1)=0 \rightarrow \tau_2 cos(\theta_2)=\tau_1 cos(\theta_1)=\tau \Rightarrow \begin{cases}
              F_y= \tau \left[tan(\theta_2)-tan(\theta_1)\right]-mg  \\
              F_x=0
            \end{cases}
          $
        }

        \textcolor{hwColor}{
          $
            F_y= \tau \left[\dfrac{\partial y}{\partial x}\Biggr|_{x=x_2}^{}-\dfrac{\partial y}{\partial x}\Biggr|_{x=x_1}^{}\right]-mg \\
            \\
            =\dfrac{\tau \left[\dfrac{\partial y}{\partial x}\Biggr|_{x=x_2}^{}-\dfrac{\partial y}{\partial x}\Biggr|_{x=x_1}^{}\right]-mg}{\Delta x}. \Delta x \\
            \\
            lim \Delta x \rightarrow 0 \Longrightarrow F_y=\left(\tau \dfrac{\partial^2 y}{\partial x^2}-mg\right).\delta x
          $
        }

        \textcolor{hwColor}{
          $F_y=m\dfrac{\partial^2 y }{\partial t^2}$ ~~~ Newton's second law and the mass of the string is $m= \mu \delta x$ \\
          \\
          By equating the two $F_y$ values we have: \\
          \\
          $
            \left(\tau \dfrac{\partial^2 y}{\partial x^2}-mg\right).\delta x=\mu \delta x\dfrac{\partial^2 y }{\partial t^2} \\
            \\
            \left(\tau \dfrac{\partial^2 y}{\partial x^2}-mg\right)=\mu \dfrac{\partial^2 y }{\partial t^2} \rightarrow \left(\tau \dfrac{\partial^2 y}{\partial x^2}-mg\right)-\mu \dfrac{\partial^2 y }{\partial t^2}=0 \\
            \\
            \\
            \dfrac{\partial^2 y}{\partial x^2}-\dfrac{\mu}{\tau} \dfrac{\partial^2 y }{\partial t^2}=\dfrac{mg}{\tau} ~~~ v=\sqrt{\dfrac{\tau}{\mu}} \\
            \\
            \Longrightarrow \dfrac{\partial^2 y}{\partial x^2}-\dfrac{1}{v^2} \dfrac{\partial^2 y }{\partial t^2}=\dfrac{mg}{\tau}
          $
        }

      \item Check if the new equation can still be solved by separation of the variables (it is not necessary to find the solution).  [ Optional: reflect on your result. Would a musical instrument sound the same on Earth and on a future human colony on the Moon?  ] 

      \end{enumerate}

    \item Generalize the problem of the vibrating string to the case of a vibrating rectangular drum of sides $a$ and $b$. Consider the 2-dimensional wave equation describing the (transverse) displacement, $U(x,y,t)$, of the drum's membrane:
    $$
    \frac{\partial^2 U}{\partial x^2} + \frac{\partial^2 U}{\partial y^2}  - \frac{1}{v^2}\frac{\partial^2 U}{\partial t^2}=0
    $$
    and start to solve it using the method of the separation of the variables. Proceed until you obtain three separate ordinary differential equations. 
    
    \item {\bf Bonus: } continue the previous problem: discuss the signs of the linking constants and show that they must take discrete values if we impose that the edges of the drum are fixed at all times. 
  \end{enumerate}

\end{document}
